
\documentclass[10pt,a4paper]{article}
\usepackage[top=12mm,bottom=12mm,left=14mm,right=14mm,headsep=2mm,footskip=7mm]{geometry}
\usepackage{fontspec}
\usepackage{xeCJK}
\setmainfont{Noto Sans}
\setsansfont{Noto Sans}
\setCJKmainfont{Noto Sans CJK SC}
\setCJKsansfont{Noto Sans CJK SC}
\usepackage{xcolor}
\usepackage{graphicx}
\usepackage{tikz}
\usetikzlibrary{calc,arrows.meta,positioning}
\usepackage{fancyhdr}
\usepackage{hyperref}
\usepackage{ragged2e}
\usepackage{microtype}
\definecolor{P2Paper}{HTML}{FCFAF7}
\definecolor{P2Ink}{HTML}{2D3238}
\definecolor{P2Muted}{HTML}{70757A}
\definecolor{P2BlueLight}{HTML}{D3EAF5}
\definecolor{P2Apricot}{HTML}{F5D9B8}
\definecolor{P2Rose}{HTML}{E8C7D3}
\definecolor{P2Violet}{HTML}{D7D3EA}
\definecolor{P2Blue}{HTML}{5B9FBE}
\definecolor{P2ApricotAccent}{HTML}{D49757}
\definecolor{P2RoseAccent}{HTML}{B97589}
\definecolor{P2VioletAccent}{HTML}{8177A8}
\definecolor{P2Rule}{HTML}{DDE3E7}
\definecolor{P2Panel}{HTML}{FFFDFC}
\pagecolor{P2Paper}
\color{P2Ink}
\setlength{\parindent}{0pt}
\setlength{\parskip}{3.5pt}
\linespread{1.04}
\hypersetup{hidelinks}
\pagestyle{fancy}
\fancyhf{}
\renewcommand{\headrulewidth}{0pt}
\fancyfoot[L]{\fontsize{6.6}{8}\selectfont\color{P2Muted}PROCESS REPORT · WORKFLOW CONSTRUCTION}
\fancyfoot[R]{\fontsize{6.6}{8}\selectfont\color{P2Muted}\thepage}
\newcommand{\eyebrow}[1]{{\fontsize{7.8}{9.4}\selectfont\bfseries\color{P2RoseAccent}#1}\par\vspace{1pt}}
\newcommand{\pagetitle}[1]{{\fontsize{18.5}{21.5}\selectfont\bfseries\color{P2Ink}#1}\par\vspace{3pt}}
\newcommand{\deck}[1]{{\fontsize{9.0}{13.0}\selectfont\color{P2Ink}\RaggedRight #1}\par\vspace{4pt}}
\newcommand{\tracehead}[2]{{\fontsize{7.1}{8.4}\selectfont\bfseries\color{#1}#2}\par}
\newcommand{\tracetitle}[1]{{\fontsize{10.2}{12.2}\selectfont\bfseries\color{P2Ink}#1}\par\vspace{1.5pt}}
\newcommand{\tracebody}[1]{{\fontsize{8.0}{11.2}\selectfont\color{P2Ink}\RaggedRight #1}\par}
\newcommand{\provenance}[1]{\vspace{2pt}{\color{P2Rule}\hrule height .35pt}\vspace{2.5pt}{\fontsize{6.9}{9.2}\selectfont\color{P2Muted}\RaggedRight\textbf{Original Evidence.} #1}\par}
\tikzset{
 userA/.style={rounded corners=1mm,draw=P2ApricotAccent,line width=.55pt,fill=P2Apricot,fill opacity=.34},
 userR/.style={rounded corners=1mm,draw=P2RoseAccent,line width=.55pt,fill=P2Rose,fill opacity=.34},
 gptB/.style={rounded corners=1mm,draw=P2Blue,line width=.55pt,fill=P2BlueLight,fill opacity=.27},
 gptV/.style={rounded corners=1mm,draw=P2VioletAccent,line width=.55pt,fill=P2Violet,fill opacity=.27},
 arrA/.style={-{Stealth[length=2.0mm,width=1.45mm]},draw=P2RoseAccent,line width=.72pt},
 arrB/.style={-{Stealth[length=2.0mm,width=1.45mm]},draw=P2Blue,line width=.72pt},
 arrV/.style={-{Stealth[length=2.0mm,width=1.45mm]},draw=P2VioletAccent,line width=.72pt},
 tagA/.style={circle,minimum size=4.8mm,inner sep=0pt,draw=P2ApricotAccent,line width=.6pt,fill=P2Paper,font=\fontsize{6.7}{7}\selectfont\bfseries},
 tagR/.style={circle,minimum size=4.8mm,inner sep=0pt,draw=P2RoseAccent,line width=.6pt,fill=P2Paper,font=\fontsize{6.7}{7}\selectfont\bfseries},
 tagB/.style={circle,minimum size=4.8mm,inner sep=0pt,draw=P2Blue,line width=.6pt,fill=P2Paper,font=\fontsize{6.7}{7}\selectfont\bfseries}
}
\begin{document}

\eyebrow{WORKFLOW CONSTRUCTION · 07}
\pagetitle{从“不确定”到候选集：AI不能替人拍板}
\deck{当 GPT 对一个研究问题并不确定时，我不希望它为了给答案而强行选一个。我的要求是：先把 2--3 个最合理的候选列出来，我们一起看依据和问题；如果讨论后仍然判断不了，再交给实验。后来为了方便记录，我们把这两种情况分别记成 CANDIDATE 和 UNRESOLVED。}
\begin{tikzpicture}[x=1mm,y=1mm]
  \node[anchor=north west,inner sep=0] (u) at (0,0) {\includegraphics[width=74mm]{assets/user_uncertainty.png}};
  \node[anchor=north west,inner sep=0] (g) at (0,-32) {\includegraphics[width=68mm]{assets/gpt_uncertainty.png}};
  \begin{scope}[shift={(u.south west)},x={($(u.south east)-(u.south west)$)},y={($(u.north west)-(u.south west)$)}]
    \path[userA] (0.31,0.33) rectangle (0.98,0.76); \coordinate (ucand) at (0.98,0.54);
  \end{scope}
  \begin{scope}[shift={(g.south west)},x={($(g.south east)-(g.south west)$)},y={($(g.north west)-(g.south west)$)}]
    \path[gptB] (0.18,0.62) rectangle (0.95,0.70); \coordinate (gdefs) at (0.95,0.66);
    \path[gptV] (0.18,0.02) rectangle (0.62,0.08); \coordinate (gcand) at (0.62,0.05);
  \end{scope}
  \draw[arrA] (ucand) .. controls (82,-3) and (82,-45) .. (gdefs);
  \draw[arrB] (ucand) .. controls (86,-7) and (86,-63) .. (gcand);
  \node[tagA] at ($(ucand)+(3.8mm,0)$) {1};
  \node[anchor=north west,text width=86mm,align=left] at (92,-2) {
    \tracehead{P2ApricotAccent}{MICRO TRACE 14}\tracetitle{“不是100\%确定，就先列2--3种”}\tracebody{我把不确定性处理从“AI继续猜”改成了“先显式给候选”。GPT随后把这一要求整理为 CANDIDATE：当多个解释都合理且会影响方法、参数或结论时，先给出候选、依据、优缺点和风险，再由我来判断。}
    \vspace{7mm}
    \tracehead{P2Blue}{HUMAN DECISION}\tracetitle{我参与的是研究判断，而不是每个操作步骤}\tracebody{这并不意味着每个小步骤都来问我。只有会改变方法、参数原则、指标或结论的问题才需要我参与；普通实现细节仍然可以自动执行。}
  };
\end{tikzpicture}
\provenance{我的原句来自真实历史截图；GPT截图来自后续整理出的不确定性处理方式。高亮与长箭头仅连接语义上直接对应的 phrase。}
\newpage
\eyebrow{WORKFLOW CONSTRUCTION · 08}
\pagetitle{讨论仍无法决定时，让实验而不是口才裁决}
\deck{我的原始要求还有后半句：如果讨论以后仍不能确认，就把不同假设并行做实验，再根据实验结果选择。这里第一次形成了后来 UNRESOLVED $\rightarrow$ 公平实验 $\rightarrow$ 数据裁决的工作流。}
\begin{tikzpicture}[x=1mm,y=1mm]
  \node[anchor=north west,inner sep=0] (u) at (0,0) {\includegraphics[width=74mm]{assets/user_uncertainty.png}};
  \node[anchor=north west,inner sep=0] (g) at (0,-32) {\includegraphics[width=68mm]{assets/gpt_parallel.png}};
  \begin{scope}[shift={(u.south west)},x={($(u.south east)-(u.south west)$)},y={($(u.north west)-(u.south west)$)}]
    \path[userR] (0.24,0.00) rectangle (0.99,0.43); \coordinate (urun) at (0.99,0.22);
  \end{scope}
  \begin{scope}[shift={(g.south west)},x={($(g.south east)-(g.south west)$)},y={($(g.north west)-(g.south west)$)}]
    \path[gptB] (0.18,0.21) rectangle (0.96,0.57); \coordinate (gpar) at (0.96,0.39);
    \path[gptV] (0.18,0.01) rectangle (0.98,0.12); \coordinate (gdata) at (0.98,0.065);
  \end{scope}
  \draw[arrA] (urun) .. controls (82,-3) and (82,-50) .. (gpar);
  \draw[arrV] (urun) .. controls (87,-8) and (87,-77) .. (gdata);
  \node[tagR] at ($(urun)+(3.8mm,0)$) {2};
  \node[anchor=north west,text width=86mm,align=left] at (92,-2) {
    \tracehead{P2RoseAccent}{MICRO TRACE 15}\tracetitle{“讨论后仍无法确认 $\rightarrow$ 并行实验”}\tracebody{GPT 后来把这种情况记成 UNRESOLVED：讨论已经不能继续提供新的判断依据时，就固定其他条件，用相同数据和指标做并行或消融实验。}
    \vspace{7mm}
    \tracehead{P2VioletAccent}{DECISION RULE}\tracetitle{方法不是由AI“硬猜”出来}\tracebody{这条规则把研究Decision的终点说清楚了：GPT可以提出候选，我负责判断哪些值得验证；如果仅靠讨论仍然不够，最后就让真实实验结果说话。}
  };
\end{tikzpicture}
\provenance{本页只使用原始要求中“讨论后仍无法确认、并行实验、根据结果选择”的部分，并连接到后续 UNRESOLVED / 实验裁决规则。}
\end{document}
